\documentclass[11pt,fontset=none]{ctexart}
\usepackage[a4paper,top=2cm,bottom=2cm,left=2.1cm,right=2.1cm]{geometry}
\usepackage{fontspec}
\usepackage{xcolor}
\usepackage{hyperref}
\setCJKmainfont[BoldFont={WenQuanYi Zen Hei}]{WenQuanYi Micro Hei}
\setCJKsansfont[BoldFont={WenQuanYi Zen Hei}]{WenQuanYi Micro Hei}
\setmainfont{DejaVu Serif}
\setsansfont{DejaVu Sans}
\hypersetup{colorlinks=true,linkcolor=blue,urlcolor=blue}
\setlength{\parindent}{0pt}
\setlength{\parskip}{0.65em}
\title{bio\_fly 教师汇报 PPT 演讲稿}
\author{bio\_fly 项目}
\date{2026 年 4 月 29 日}
\begin{document}
\maketitle
本文档对应 `/unify/ydchen/unidit/bio\_fly/ppt/cyber\_fly\_teacher\_report.tex`，按页给出汇报口径。整体讲法是：先讲问题，再讲结构证据，再讲连接组传播，再讲行为框架，最后主动说明边界和后续验证。
\bigskip\hrule\bigskip
\section*{第1页 封面}
今天我汇报的不是一个单纯的“赛博果蝇视频工程”，而是一条从连接组统计、功能传播，到嗅觉记忆行为框架的完整证据链。核心问题只有一个：左右脑蘑菇体的递质输入侧化，能不能进一步影响果蝇记忆相关的功能读出。
如果把这项工作用工程语言概括，就是我们把公开 FlyWire 数据、稀疏传播、行为代理和视频渲染串成了一个可复算的 pipeline；如果用生物学语言概括，就是我们在检验“结构侧化是否具有功能可传播性”。
\section*{第2页 先把名词说清楚 1：神经递质和侧化}
这页的目的，是先把“我们到底在量什么”说清楚。这里的神经递质不是泛泛的脑内化学物质，而是指 KC 上游输入里各类递质的预测占比。
“侧化”不是说左右脑大小不同，而是说同一个脑区在左右两侧接收到的输入化学组成不同。我们今天最重要的两个方向是：右侧血清素更高，左侧谷氨酸更高。
\section*{第3页 先把名词说清楚 2：递质名称不要混淆}
这页主要是避免术语误解。Ser、5-HT 指的是血清素，Glu、GLU 指的是谷氨酸；DA 是多巴胺，GABA 是抑制性递质，ACh 是乙酰胆碱。
特别要强调的是，OCT 气味是 3-octanol，MCH 气味是 4-methylcyclohexanol，它们是行为实验中的气味刺激，不是章鱼胺，不要和 Oct/OA 混在一起。
\section*{第4页 先把名词说清楚 3：脑区和统计量}
这一页是为了让不做生物的人也能跟上。MB 是蘑菇体，KC 是它的主细胞，MBON 是输出，DAN 是多巴胺教学信号，APL 和 DPM 更偏反馈和记忆维持，DN 是下降神经元，也就是从脑到身体动作系统的重要出口。
统计量上，Cohen's d 用来表示左右差异的效应量，FDR q 是多重比较校正后的显著性，laterality index 就是左右偏向指数。我们后面所有结论，都会尽量落在这些可复算的量上。
\section*{第5页 本研究最核心的原始发现}
这一页先给出最核心的主结果：蘑菇体 KC 的上游输入不是左右完全对称的。最稳定的方向是右侧血清素富集，左侧谷氨酸偏置，而且这个模式不是只出现在单个细胞类型里，而是跨 KC 亚型反复出现。
从工程上讲，这意味着输入特征本身已经带有方向性；从生物学上讲，这意味着蘑菇体的左右半球不是简单复制件，而可能存在输入化学组成上的分工。
\section*{第6页 结构证据：血清素右偏、谷氨酸左偏不是孤立个例}
这一页的重点是说明：这个现象不是偶然点状出现，而是一个整列模式。图里红色代表右偏，蓝色代表左偏，我们看到 Ser/5-HT 基本整列偏右，Glu/GLU 多数偏左。
所以我们不是在讲“某一个亚型碰巧有差异”，而是在讲“整个 KC 输入系统存在稳定的左右偏向”。这也是为什么后面可以把它拿来做功能传播，而不是停在描述性统计上。
\section*{第7页 为什么这和记忆有关}
这页要把结构发现和记忆生物学连起来。为什么我们关心 $\alpha'\beta'$？因为它本来就和记忆巩固、稳定化相关；而我们的侧化信号在这个亚区最强。
这意味着这不是单纯的侧化热图，而是一个和记忆过程有关的侧化。换句话说，它给了我们一个可测试的假说：左右半球可能不是完全冗余，而是对记忆处理承担不同的输入化学角色。
\section*{第8页 仿真如何检验该发现的功能意义}
这里开始进入连接组传播。我们的做法不是随机激活一大片蘑菇体，而是选结构上有依据的 seed：右侧血清素富集 KC 和左侧谷氨酸富集 KC。
然后把这些 seed 放进 FlyWire 的 signed connectivity 上做多跳稀疏传播，再和同侧、同规模随机 KC 对照比较。这样做的目的，是区分“递质侧化效应”与“一般性 KC 激活效应”。
\section*{第9页 仿真结果：递质侧化能传播到记忆相关读出}
这一页是功能层结果。右侧血清素 KC 激活后，DAN、DPM、APL、左右偏向这些读出都会变化；左侧谷氨酸 KC 激活后，memory axis、DAN、MBIN、DPM、MBON、APL 和左右读出都会变化。
这说明侧化不是停留在输入端，而是可以沿连接组传播到下游功能读出。对工程师来说，这是一个从输入特征到下游响应的可解释映射；对生物学来说，这是侧化具有功能可传播性的证据。
\section*{第10页 这轮仿真的支持范围与边界}
这一页必须主动讲清楚边界。我们能支持的是：结构侧化稳定存在，功能读出上也能传播；但我们不能说已经证明真实行为因果。
同样，公开资料没有完整 DN-to-body 权重，所以当前实现的是公开可审计的 surrogate，而不是 Eon 私有闭环的逐参数复刻。这个边界讲清楚，整篇故事反而更稳。
\section*{第11页 当前最稳的结论}
这一页可以作为汇报中的“结论摘要”。第一，结构层上，血清素右偏、谷氨酸左偏最强出现在 $\alpha'\beta'$；第二，功能层上，这种侧化可以传播到 memory axis 和 MB family 相关读出；第三，行为层上，OCT/MCH 任务本身是跑通的；第四，边界上，MB 扰动目前还没有显著行为效应。
所以现在最稳的不是“我们已经证明一切”，而是“我们已经建立了一条结构到功能、再到行为框架的可证伪证据链”。
\section*{第12页 研究主线与证据链}
这页我建议直接按证据链讲。第一层是结构不对称；第二层是连接组传播；第三层是行为任务；第四层是因果边界。
对老师来说，这种讲法比直接说“赛博果蝇自动涌现行为”更可信；对工程师来说，这种讲法也更符合系统调试思路，因为它把每一层的输入、输出和失败模式都拆开了。
\section*{第13页 关键变量定义}
这一页是给第一次接触这个项目的人看的字典。最关键的是 MB、KC、MBON、DAN、APL/DPM、DN 这些缩写的角色分工。
如果只记一句话：KC 负责把气味变成稀疏编码，MBON/DAN/APL/DPM 负责价值和记忆调节，DN 是从脑到身体行为的出口。后面所有实验，都是围绕这条链路展开的。
\section*{第14页 本研究包含的实验类型}
这一页告诉大家我们不是只做了一张图，而是做了四类实验：结构侧化审计、四卡连接组传播、OCT/MCH 行为仿真、MB 侧化扰动，还有视频和图表整理。
从系统角度看，这些模块各自回答不同问题：结构是否不对称、功能是否传播、行为任务是否成立、扰动是否能打出因果差异。这样搭起来，论文结构就完整了。
\section*{第15页 数据和算力：这不是手工玩具模型}
这一页主要是让工程背景的老师知道，这不是几张图拼起来的展示，而是一套有数据规模和算力支撑的流程。FlyWire 神经元数量、连接边数、MB family 的规模都已经足够大，足以支持统计和传播实验。
我们还有四张 GPU 可并行跑 trial sweep 和随机对照，这保证了不是单次样本碰运气。但同时也要说明，这还不是 Eon 那种私有闭环控制器的逐参数复刻。
\section*{第16页 实验 1：先看结构不对称是否真实存在}
这一页进入第一个正式实验。方法很简单：按 MB family 分组，统计左右连接强度，再算 right laterality index。
这一步的作用是先回答最基础的问题：左右蘑菇体到底是不是完全镜像。只有这个问题先站住，后面的传播和行为讨论才有基础。
\section*{第17页 实验 1 的发现：结构不是完全镜像}
这里的结果很直接：KC 内部复发、KC 到 DPM、APL 到 KC、DPM 到 KC 多数呈左偏，而 MBON 到 DAN、DAN 到 MBON、MBON 到 APL 出现轻度右偏。
所以可以把结构层粗略理解成：左侧更像反馈/维持轴，右侧更像输出/教学轴候选。但这还只是结构推断，不能直接跳成行为结论。
\section*{第18页 实验 2：把 MB 记忆轴传播到 DN}
第二个实验开始把结构和功能接起来。输入 seed 分成 MBON、DAN、APL/DPM 和 memory axis，再加一个 OCT/MCH KC context 作为气味上下文。
然后通过 FlyWire signed connectivity 做 3 hop propagation，并在 DN 层读取响应量、总质量和左右偏向。这个实验的目标就是看 MB 的记忆轴能不能通向动作出口。
\section*{第19页 实验 2 的发现：MBON 是最强 DN 出口}
这一页最重要的是区分两个概念：输出强度和左右偏向。left MBON 的总 DN 响应最强，但 right MBON 和 right memory axis 的右偏更明显。
所以我们不能只看“谁总量最大”，还要看“谁更偏向哪一侧”。这也是为什么 laterality index 在这里非常重要，它帮我们把强度和方向拆开了。
\section*{第20页 实验 2 的机制图：为什么要加 DN/motor primitive}
这页的逻辑是：公开资料没有给完整 DN-to-body 权重，所以我们不能冒充自己已经恢复了一个私有闭环控制器。
因此我们采用透明的中间层，把 MB seed 映射到 DN family，再映射到 motor primitive。这样每一步都能复算，也便于以后用真实行为或成像数据替换。
\section*{第21页 实验 2 的行为出口：现在只到低维动作假说}
这一页要主动降调。我们现在能说的是：MBON/memory axis 主要读出到 DNa/DNp 等下降通路，当前最稳定的 motor primitive 是 turning / orientation 候选。
但不能把这张热图直接当成真实果蝇肌肉命令，更不能把它说成完整行为涌现。它现在更像一个可解释的动作低维读出层。
\section*{第22页 实验 3：OCT/MCH 嗅觉记忆怎么设计}
第三个实验是行为框架。这里要先解释 CS+、CS-、US：OCT+sucrose 代表气味和糖奖励配对，OCT 是 CS+；OCT+shock 代表气味和电击配对，OCT 仍然是 CS+，但要回避。
我们还设计了 mirror-side，把 CS+ 左右位置互换，避免把位置偏好误当成气味记忆。对工程背景的人来说，这就是在控制一个很重要的混杂变量。
\section*{第23页 实验 3 的发现：奖励和惩罚方向显著}
这一页说明行为框架本身是有效的。奖励条件会接近 CS+，惩罚条件会回避 CS+，冲突条件仍然能表现出记忆驱动的选择。
这意味着我们的行为代理至少能表达 valence，也能表达基本记忆方向；所以它适合做方法展示和补充材料，但它不是最终因果证据。
\section*{第24页 最重要的负结果：MB 扰动还没有打出显著行为差异}
这一页很关键，必须主动讲。当前 left/right silencing、symmetrization、swap 等 MB 扰动，在行为层没有打出显著差异。
这不是坏事，它实际上限定了故事边界：说明现有行为代理对侧化 readout 的灵敏度还不够，不能把结构差异直接说成已被行为因果证明。负结果写清楚，论文反而更可信。
\section*{第25页 会议反馈后的新增实验：5-HT 与 Glu 分拆验证}
这一页是根据最近一次生物老师反馈新加的。老师们指出，不能把“右侧血清素”和“左侧谷氨酸”混成一个模糊的侧化变量，而应该拆成两条可验证机制轴。
所以我们重新按绝对效应量 $|z|$ 做 readout 比较。这里要解释清楚：$z$ 的正负表示读出方向，不表示谁更重要；比较强弱时看绝对值。结果显示，左侧谷氨酸是更强的广谱 memory-output 扰动，而右侧血清素更适合进入 DPM、5-HT 释放和记忆巩固时间窗验证。
\section*{第26页 会议反馈后的新增实验：DPM 光遗传协议扫描}
这一页回应的是生物老师提出的一个具体实验：既然成像后果蝇不能再做行为，那就先在同一只果蝇上通过光遗传反复刺激 DPM，看 5-HT 或下游响应模式是否稳定。
我们把这个建议转成了仿真实验：用 GPU0 和 GPU1 分别传播左侧和右侧 DPM seed，再模拟不同频率、时长和波形的刺激。结果很清楚，left DPM 和 right DPM 的 readout laterality 方向相反，这给后续成像实验提供了可检验预测。
\section*{第27页 DPM 光遗传验证：先证明 release pattern}
这一页要讲清楚为什么我们把成像证明和行为证明分开。脑成像会破坏或强扰动果蝇，所以不能要求同一只果蝇先测 5-HT 释放再继续做行为。
我们建议先做 DPM-driver 驱动 CsChrimson 或 ReaChR，在 KC 或蘑菇体 compartment 表达 5-HT sensor 或 GCaMP。主指标不是看一张漂亮图，而是定量左右 release LI、peak dF/F 和 AUC。最关键的控制是把果蝇水平旋转 180 度后按脑侧注册，如果是真实偏侧化，脑侧符号应保持；如果是成像角度伪影，图像坐标符号会翻转。
\section*{第28页 DPM 光遗传验证：再做群体行为证明}
这一页讲第二条证明链。行为不和破坏性成像绑定，而是在独立果蝇群体中做 OCT/MCH T-maze 或相机轨迹实验。
仿真预测最敏感的是 weak OCT / strong MCH conflict 条件，因为普通 OCT/MCH choice rate 容易饱和。用 617 或 627 nm 红光刺激 DPM，预测 choice-index delta 大约是 +0.10。这个数值不是已完成湿实验结果，而是给老师设计样本量和条件优先级的效应方向。
\section*{第29页 DPM 光遗传验证：推荐协议和实验链}
这一页给出最直接可操作的协议：优先测试 CsChrimson 617 nm 或 ReaChR 627 nm，40 Hz，20 ms pulse，5 s，0.1 mW/mm2。蓝光 ChR2 可以作为 positive control 或校准，但不建议作为主行为实验，因为蓝光更容易引入视觉混杂。
这里的核心逻辑是三段式：GRASP 或 split-GFP 做结构硬证据，DPM 光遗传成像做 release pattern 功能证据，独立群体 T-maze 做行为调节证据。三者合在一起，比单独追求“同一只果蝇全都做完”更现实，也更严谨。
从第30页开始，原讲稿页码相对旧版顺延 3 页；汇报时以新版 42 页 PDF 的页码为准。
\section*{第27页 会议反馈后的新增实验：群体 T-maze 可观测指标}
这一页解决另一个现实问题：单只脑成像和几百只果蝇的行为统计规模对不上，所以不能做“同一只果蝇先成像再行为”的关联。
我们的折中方案是把仿真结果压缩成群体 T-maze 可测的 choice index。预测是：野生型 choice index 最高；单独削弱 5-HT-right 或 Glu-left 都会下降；两条轴一起削弱下降最明显。这个结果不能单独证明偏侧化，但可以指导真实群体实验该看什么指标、预期差异多大。
\section*{第28页 会议反馈后的验证逻辑：结构是硬红线}
这一页是整套故事中最重要的证据分级。结构层的 GRASP 或 split-GFP 是硬证据；DPM 光遗传、钙成像和 T-maze 行为是功能和行为佐证。
因此我们现在的严谨说法是：仿真把结构发现转化成可操作实验预测，但不能替代结构验证。这个口径非常重要，因为它避免把连接组模型讲成已经完成真实生物因果证明。
\section*{第29页 视频做了什么：让实验能被看懂，而不是替代统计}
这一页是给老师看的“视频用途说明”。视频补了培养皿、气味杯、滤纸、气味羽流、糖滴、电击栅格和轨迹尾迹，让实验场景一眼能看懂。
但必须明确，视频只是展示和检查，不是显著性证据。真正的统计结论仍然来自 trial 级别的正式分析。
\section*{第30页 把三层证据连起来}
这一页是整套故事的结构图。结构层说明左右 MB 不对称，功能层说明这种不对称能传播到 DN，行为层说明 OCT/MCH 任务本身是可工作的。
这张图要强调的不是“完美闭环”，而是“证据链”。我们现在已经能把结构、传播和行为串起来，但因果层还需要更强的实验支撑。
\section*{第31页 目前真正的新发现是什么}
这一页可以概括为四个层次：递质侧化、结构连接不对称、功能传播、行为框架。真正的新发现不是单点，而是这些层次能被串起来。
同时也要主动说约束性发现：当前行为代理对 MB 侧化扰动不敏感，所以还不能把递质侧化说成已经被行为因果证明。这个边界很重要。
\section*{第32页 当前最稳的结论层级}
这一页我建议作为汇报中的“最终稳结论”。结构层、功能层、边界层、方法层分别对应不同证据强度。
如果只让听众记一句话，就是：我们已经证明了一个稳定的结构侧化，并且证明它可以传播到功能读出；但行为因果还需要更强的实验闭环来确认。
\section*{第33页 这项工作的意义}
从神经生物学角度，这个工作提出了一个具体的、可检验的假说：右侧血清素调节和左侧谷氨酸输入，可能给蘑菇体记忆环路提供左右分工。
从计算建模角度，我们把“递质输入组成不对称”转成了可运行、可统计、可消融的连接组传播实验。对论文叙事来说，这比只讲一个结构热图更完整，也更接近探索性机制论文的写法。
\section*{第34页 主要限制与待补证据}
这一页要主动说明系统限制。第一，公开资料没有完整 DN-to-body 权重，所以我们不能声称复刻了私有闭环；第二，视频和轨迹仍然带代理层，不是全生物真实模型；第三，递质预测本身有误差。
所以当前最合理的表述是：我们已经有了一个公开可审计的 surrogate，并且它能支持假说生成，但还不能替代湿实验。
\section*{第35页 后续计算工作}
如果继续做计算层，我认为最值得优先补的是三件事：第一，把 MB-DN readout 接入控制器，让它不只是解释层；第二，继续加强 MB 扰动对照，做更严格的随机化和对称化；第三，扩展行为背景随机性，检验稳健性。
这三件事的目标不是“再做大一点”，而是提高因果解释的分辨率，让模型能更接近可检验的闭环。
\section*{第36页 后续湿实验验证}
如果要把这条故事推进到更强的论文等级，最终还是要回到湿实验。最关键的是 GRASP 结构验证：优先看 right DPM/5-HT input 到 right KCa'b'，并用 left Glu input 到 left KCa'b' 作为相反方向 positive control。
功能层则优先做 DPM 光遗传协议扫描、左右 KC/MBON/DN 钙成像，以及 OCT/MCH T-maze 群体 choice index。只有结构、功能和行为三层对齐，侧化与记忆的因果链才算站稳。
\section*{第37页 汇报结构建议}
这一页其实是给我自己和大家一个讲法建议。先讲科学问题，再讲结构侧化，再讲 MB-DN readout，再讲 OCT/MCH 行为，最后主动说 MB 扰动没有显著，并补充会议反馈后的 GRASP、DPM 光遗传和 T-maze 预测。
这样讲的好处是证据强弱顺序清楚，而且不会让听众误以为我们在把一个代理系统包装成真实因果系统。对工程师和 AI 背景的老师，这种分层非常重要。
\section*{第38页 交付物在哪里}
这一页只是把最重要的文件和路径摆出来，方便后续查阅和复现。包括 README、教师版说明、实现状态报告、MB-DN 报告、会议反馈报告、论文草稿、英文 PDF、中文解读 PDF、视频和 PPT。
如果老师或工程师后续要复跑，最关键的是看这些绝对路径，而不是看截图。这样能够减少“资料散落”的问题。
\section*{第39页 结论与展望}
最后我会用两句话收尾。第一，我们现在已经建立了一个公开可复现的 FlyWire 连接组到行为框架，它支持“右侧血清素富集、左侧谷氨酸偏置的递质侧化可以传播到记忆轴和 DN 出口”的新假说。
第二，我们还不能说已经完整复刻私有赛博果蝇，更不能说已经证明连接组单独自动涌现完整行为。下一步最关键的问题，是用 GRASP 验证结构侧化，并用 DPM 光遗传、钙成像和群体 OCT/MCH 行为去检验这些预测。
\end{document}
